% Generated draft: requires mathematical and source review.
\documentclass[11pt]{article}
\usepackage[margin=25mm]{geometry}
\usepackage{fontspec}
\setmainfont{Latin Modern Roman}
\usepackage{amsmath,amssymb,mathtools,unicode-math}
\setmathfont{Latin Modern Math}
\usepackage{enumitem,needspace,xurl,hyperref}
\setlength{\parindent}{0pt}
\setlength{\parskip}{5pt}
\setlength{\emergencystretch}{3em}
\newcommand{\N}{\mathbb{N}}
\newcommand{\Z}{\mathbb{Z}}
\newcommand{\Q}{\mathbb{Q}}
\newcommand{\R}{\mathbb{R}}
\newcommand{\C}{\mathbb{C}}
\newcommand{\F}{\mathbb{F}}
\newcommand{\A}{\mathbb{A}}
\newcommand{\PP}{\mathbb{P}}
\newcommand{\E}{\mathbb{E}}
\newcommand{\eps}{\varepsilon}
\newcommand{\dd}{\,\mathrm{d}}
\begin{document}
\section*{UnsolvedMath ID: 700438}
\textbf{Source problem code:} catalogue entry
\subsection{Definitions and mathematical statement}
A group \(G\) is linear over a field \(K\) if there are an integer \(n\geq 1\) and an injective homomorphism \(G\to\mathrm{GL}_n(K)\), where \(\mathrm{GL}_n(K)\) is the group of invertible \(n\times n\) matrices over \(K\). The supplied record does not specify a field or characteristic restriction.

Let \(\mathcal C_I\) and \(\mathcal C_N\) denote, respectively, the classes called \(\overline{\!SI\vphantom{(}}\)-groups and \(\overline{\!SN\vphantom{(}}\)-groups in the source. Their defining properties, including the meaning of the overlines, are not supplied; these symbols therefore remain unresolved source terminology.

\textbf{Question (source problem 7.41).} Under the usual interpretation of linearity over arbitrary fields, is it true that, for every field \(K\), every integer \(n\geq1\), and every subgroup \(G\leq\mathrm{GL}_n(K)\),
\[
G\in\mathcal C_I\quad\Longrightarrow\quad G\in\mathcal C_N\,?
\]
This is a schematic restatement: a fully defined mathematical assertion requires the source definitions of both group classes and confirmation of the intended convention for linearity.
\subsection{Short English statement}
The question asks whether every linear \(\overline{\!SI\vphantom{(}}\)-group is an \(\overline{\!SN\vphantom{(}}\)-group. The supplied record does not define these two classes, so their precise meaning must be recovered before this draft can receive mathematical review.
\subsection{Sources}
\begin{itemize}
\item Supplied local UnsolvedMath registry record: permanent ID \(700438\), exact source code \texttt{catalogue entry}, title \emph{Kourovka Notebook Problem 7.41}; upstream revision \texttt{d371ee917ac9647fbb06cd6fc529deec6f3dda82}. The record labels the problem open; this is supplied registry metadata, not an independently verified mathematical assessment.
\item Supplied bibliographic attribution: \emph{The Kourovka Notebook (Unsolved Problems in Group Theory)}, 7th issue (1980), problem 7.41; also identified with the 21st issue edition (2026), arXiv:1401.0300v48. The collection description names E. I. Khukhro and V. D. Mazurov as editors.
\item The background and proposer field attribute the problem to Yu. I. Merzlyakov, with proposed year 1980. The statement itself begins with the attribution ``(John S. Wilson).'' The supplied record does not explain the relationship between these attributions.
\end{itemize}
\end{document}
